\section{Use shortcuts but double check (lossless)：speculative sampling}

Speculative sampling 利用一个便宜 draft model \(p\) 一次猜多个 token，再让 target model \(q\) 并行检查。直觉是：generation 慢，但 checking 多个候选 token 比逐个生成便宜。关键性质是它仍然精确采样自 \(q\)，因此属于 lossless shortcut。

\begin{knowledgebox}{课堂提示：checking 比 generation 更容易并行}
老师先回顾 prefill 可并行且 compute-bound、generation 逐 token 且 memory-bound，再得出不对称性：checking 一串候选比逐个生成更快。Draft model 只负责便宜地提案，target model 仍决定分布；修改过的 rejection sampling 保证最终是 target 的 exact sample，因此这里的 speedup 不靠牺牲输出分布。
\end{knowledgebox}

\begin{figure}[H]
\centering
\includegraphics[width=0.78\textwidth]{speculative-sampling-algorithm.png}
\caption{Speculative sampling 算法：draft model 提案，target model 校验，必要时从 residual distribution 采样。}
\figsource{CS336 Lecture 10 官方讲义引用图。}
\end{figure}

\begin{knowledgebox}{读图：为什么它仍是 exact sample}
若 draft \(p\) 对某 token 给出概率高于 target \(q\)，该 token 只以 \(q/p\) 的概率接受；被拒绝的概率质量会转移到 residual distribution \(\max(q-p,0)\)。这样最终每个 token 的边际概率仍等于 \(q\)，只是把多个候选的校验并行化。
\end{knowledgebox}

\[
P(A)
  =
p(A)\frac{q(A)}{p(A)}
  +
p(B)\cdot 1\cdot 0
  =
q(A).
\]
这个两 token 例子说明：如果 \(p(A)>q(A)\)，过多提出的 \(A\) 会被拒绝一部分；若 \(p(B)<q(B)\)，不足的概率通过 residual 补回来。

\begin{figure}[H]
\centering
\includegraphics[width=0.9\textwidth]{speculative-sampling-results.png}
\caption{Speculative sampling 结果：draft/target 组合可显著减少 target model 的 sequential generation 次数。}
\figsource{CS336 Lecture 10 官方讲义引用图。}
\end{figure}

结果图中的 speedup 需要拆成三项：draft 生成候选的成本、target 并行验证候选的成本，以及一次能接受多少 token。若 draft 太弱，接受率低；若 draft 太大，提案本身昂贵；若 batch 或序列太短，target verification 的并行优势也无法摊薄。

\begin{figure}[H]
\centering
\includegraphics[width=0.9\textwidth]{speculative-sampling-stats.png}
\caption{Speculative sampling 统计：性能取决于 draft model 接近 target 的程度和候选长度。}
\figsource{CS336 Lecture 10 官方讲义引用图。}
\end{figure}

\begin{knowledgebox}{读图：speculative speedup 的两个控制旋钮}
Draft model 越接近 target，接受率越高；一次猜的 tokens 越多，潜在并行度越高，但被拒绝时浪费也越多。最佳点取决于 draft 成本、target 成本、接受率和服务端 batch 调度。
\end{knowledgebox}

\begin{figure}[H]
\centering
\includegraphics[width=0.85\textwidth]{medusa-eagle.png}
\caption{Medusa 和 EAGLE：改进 draft model 的方法，让提案更接近 target。}
\figsource{CS336 Lecture 10 官方讲义引用图。}
\end{figure}

\begin{knowledgebox}{读图：Medusa/EAGLE 的共同目标}
Medusa 让 draft 头并行提出多个未来 token；EAGLE 利用 target model 的高级特征帮助 draft。二者都服务同一个目标：提高接受率，同时保持 draft 成本足够低。
\end{knowledgebox}

\subsection{本章小结}

Speculative sampling 的美感在于：它利用系统不对称性，checking 比 generation 更可并行；又用 rejection-sampling 逻辑保证目标分布不变。它是本讲中最典型的“系统技巧 + 概率正确性”结合。

